\markboth{DATA SOURCES AND INTEGRATION}{DATA SOURCES AND INTEGRATION}
\section{Data Sources and Integration}\label{sec:data}

Training Copilot integrates eight data sources through the uniform MCP interface described in Section~\ref{sec:architecture:mcp}. Each source is connected through an independent server responsible for authentication, rate limiting, error recovery, and data normalisation. Table~\ref{tab:data_summary} summarises the sources, their tool counts and their main metrics. With one exception, all servers were implemented as native FastMCP services. The Strava server uses the official v3 REST API, while the Garmin server relies on the open-source \texttt{garminconnect} library. Weather data is obtained from Open-Meteo, and route generation combines OpenRouteService with OSM Overpass. The Google Maps Platform provides place, directions, and geocoding data, while calendar information is accessed through the Google Calendar API. The flythrough render server (Section~\ref{sec:architecture:supporting}) is likewise ours. The Telegram bridge described in Section~\ref{sec:data:telegram} is the only third-party server. We vendored the existing MCP project without modification and connected it through our transport layer. Its inclusion demonstrates that the architecture can integrate an external server in the same way as a native service. Every server runs as an independent process over Streamable HTTP and, in the containerised deployment, as its own Docker container (one per MCP server, alongside one per agent) built from a single shared Dockerfile (Section~\ref{sec:architecture:deployment}).\footnote{For all development and evaluation we used the authors' own Strava and Garmin accounts and personal training data; access to these accounts can be provided to graders on request.} This section discusses the integration requirements of each source. Table~\ref{tab:api_overview} in Appendix~\ref{sec:appendix:apis} adds a per-API reference of access requirements (API keys, OAuth, credentials), cost and rate limits, and native MCP availability for every external service used, including the LLM backend.

\subsection{Strava: Activity Tracking}\label{sec:data:strava}

The Strava server (port 8103) connects to the Strava v3 REST API via the standard OAuth\,2.0 Authorization Code flow. On first use, it opens Strava's consent page, where the user grants the required scopes (\texttt{read}, \texttt{activity:read\_all}, \texttt{activity:write}). The server then exchanges the authorization code for access and refresh tokens, stores them in \texttt{.tokens/strava.json}, then automatically refreshes them five minutes before expiry. This flow (\texttt{auth/strava\_oauth.py}) is transparent to the MCP tools, which simply call the API with a valid token.

Its thirteen tools cover recent activities, aggregate statistics, year-over-year breakdowns, personal records, gear mileage, and detailed activity data, including lap splits, power metrics, and sub-second GPS streams. They also provide performance-trend analysis against personal baselines and calculate Training Stress Balance metrics, including CTL, ATL, and TSB, using the fitness–fatigue impulse-response model~\citep{morton_modeling_1990}. HTTP errors use exponential backoff for server errors (5xx) but return immediately on rate limits (429), since Strava's \texttt{Retry-After} is typically 15\,+ minutes, which makes in-request retries pointless. In June 2026, Strava announced that API access, previously free, now requires a subscription for Standard Tier developers, alongside an official Strava MCP server for AI-native access~\citep{strava_api_update_2026}. This fits our MCP-first architecture, since the official server could eventually replace ours with a registry URL change and no agent or UI code changes. Figure~\ref{fig:screenshot_activity} in Appendix~\ref{sec:appendix:screenshots} shows the resulting Dashboard tab, combining an activity map, recent-activity list, and training-load trend charts. A dedicated \textbf{Activity Analysis} view renders per-stream charts (heart rate, pace, elevation, cadence, power) alongside a route map that recolours by the selected metric. A companion \textbf{Sync tab} (Figure~\ref{fig:screenshot_sync}) transfers activities from Garmin to Strava without manual file downloads. It previews Garmin activities for a date range (each a card with its route map and a flag for whether it is already on Strava) and uploads the selected FIT files directly to Strava, for athletes who record on Garmin but keep their history on Strava.

\subsection{Garmin Connect: Wellness and Recovery}\label{sec:data:garmin}

Garmin provides no public REST API for the wellness data (sleep, HRV, Body Battery, stress) Training Copilot needs, since the Garmin Health API is enterprise-only, so our server relies on the open-source \texttt{garminconnect} Python library,\footnote{\url{https://github.com/cyberjunky/python-garminconnect}} which reverse-engineers the Garmin Connect web endpoints. Authentication uses Garmin's SSO login with stored credentials and MFA on first login. The session tokens are cached and reused until expiry. This carries inherent risk, as the endpoints are undocumented and may change without notice. Rate limits are less predictable than Strava's (the Body Battery endpoint rejects ranges over 28 days, so the server chunks requests automatically), and exponential-backoff retries handle transient failures. A mock-data mode (\texttt{GARMIN\_MOCK\_HEALTH=true}) enables UI development without a device.

Its thirteen tools cover seven categories, including sleep, SpO\textsubscript{2}, and HRV metrics with personal baselines and readiness status, Body Battery, stress and heart-rate timelines, VO\textsubscript{2}max, training load, race predictions, body composition, and multi-day wellness trends with step and GPS data. These metrics let the recovery specialist assess training readiness by considering sleep quality, HRV relative to baseline, stress, and remaining energy together, the multi-signal approach \cite{rothschild_predicting_2024} showed is needed for accurate recovery prediction. Figure~\ref{fig:screenshot_health} in Appendix~\ref{sec:appendix:screenshots} shows the resulting Health tab.

\subsection{Weather, Routes, Calendar, and Maps}\label{sec:data:other}

The \textbf{weather server} (port 8101) queries the Open-Meteo API, which needs no API key and imposes no rate limits for non-commercial use. Its four tools provide current conditions, multi-day forecasts (1-16 days, hourly or daily) with temperature, precipitation probability, wind speed, and UV index, pollen concentrations by species, and UV category classification. The context specialist applies simple rules to classify training windows: 5-20\,\textdegree C as ideal, precipitation probability above 30\,\% as a rain warning, and UV index above 6 as a sun-protection advisory.

The \textbf{routes server} (port 8102) aggregates three external services across seven tools. OpenRouteService (via \texttt{ORS\_API\_KEY}) provides routing for cycling, running, and hiking, along with circular routes, elevation profiles, and isochrone maps. Google Geocoding converts place names to coordinates, while OSM Overpass retrieves boundaries for location-constrained routes, such as a loop within Schlossarten (a ``run inside Schlossgarten'' query generates the park boundary from OpenStreetMap, then routes within it). Each tool returns GeoJSON-compatible coordinate arrays rendered as interactive Folium maps; Figure~\ref{fig:screenshot_routes} in Appendix~\ref{sec:appendix:screenshots} shows the resulting Routes tab.

The \textbf{calendar server} (port 8105) provides six tools for Google Calendar read/write access, supporting both a single-user development mode (persisted OAuth tokens in \texttt{.tokens/google.json}) and a multi-tenant mode (per-request \texttt{Authorization: Bearer} headers). The context specialist cross-references free calendar windows with weather forecasts to suggest optimal training times and can create events for planned sessions.

The \textbf{Google Maps server} (port 8108) adds six tools over the Google Maps Platform (Places, Routes, and Geocoding APIs, all usable with a billing-free demo key): free-text place search, place details, forward and reverse geocoding, point-to-point directions with travel time, and, most useful for us, \texttt{maps\_search\_along\_route}, which finds places (a caf\'e, a bakery, a drinking-water fountain) that lie \textit{on} a planned route rather than elsewhere in town. The route tools emit a short list of \texttt{poi\_anchors} (real track points with their kilometre mark). The route specialist passes these to the search, and only places within a set detour of the track come back, so ``plan a 40\,km loop with a caf\'e stop on the way'' is answered from the actual route. Like every server except the Telegram bridge, this is our own native FastMCP code. It replaces an archived third-party proxy but keeps its tool names, so the route specialist's prompt is unchanged.

\subsection{Fitness Literature: RAG Knowledge Base}\label{sec:data:rag}

The fitness specialist (port 9005) is the only agent without an MCP server: it queries a local vector database of fifty public-domain books on physical culture, athletics, physiology and health from Project Gutenberg\footnote{\url{https://www.gutenberg.org/}} via RAG~\citep{lewis_retrieval-augmented_2020}. Embeddings come from the local \texttt{all-MiniLM-L6-v2} model ($\sim$90\,MB)~\citep{reimers_sentence-bert_2019}; the corpus is split offline into $\sim$900-character chunks with 150-character overlap ($\sim$32k chunks in total), and each query retrieves its top five chunks by cosine similarity, a single NumPy matrix-vector product that needs no FAISS index or vector-database service at this size. The specialist grounds every claim in the retrieved passages and cites the source book; an \texttt{ensure\_sources()} step keeps those citations from being dropped when the orchestrator synthesises the answer. We do this because a general-purpose model asked for exercise-science advice from memory readily produces fluent but unspecific or invented guidance~\citep{puce_harnessing_2025}; tying each recommendation to a real, citable passage keeps it traceable and surfaces an unsupported claim instead of hiding it. This is the same stance as the rest of the system, where every data answer comes from a live tool call, and an agent that cannot retrieve what it needs returns an explicit error rather than inventing one (Section~\ref{sec:agents:observability}). Whether that grounding holds is checked in part today by the deterministic \texttt{grounded\_in\_real\_data} scorer (Section~\ref{sec:evaluation:scorers}); a citation-faithfulness scorer and a dedicated answer-verification pass are left as outlook (Section~\ref{sec:conclusion:next}). Figure~\ref{fig:rag_pipeline} illustrates the pipeline.

\begin{figure}[ht]
\centering
\resizebox{\textwidth}{!}{%
\begin{tikzpicture}[
    node distance=0.4cm and 0.5cm,
    every node/.style={font=\small},
    block/.style={draw, thick, rounded corners=3pt, minimum height=0.7cm, minimum width=2.0cm, align=center},
    data/.style={draw, thick, rounded corners=1pt, minimum height=0.5cm, align=center, font=\scriptsize},
    arr/.style={-{Stealth[length=5pt]}, thick},
]
    \node[block, draw=tcblue, fill=tcblue!8] (query) {User Query};
    \node[block, draw=tcgreen, fill=tcgreen!8, right=1.0cm of query] (embed) {Embed\\{\scriptsize all-MiniLM-L6-v2}};
    \node[block, draw=tccyan, fill=tccyan!8, right=1.0cm of embed] (store) {Cosine Search\\{\scriptsize (N$\times$384 float32)}};
    \node[block, draw=tcyellow!80!black, fill=tcyellow!8, right=1.0cm of store] (topk) {Top-$k$\\Passages};
    \node[block, draw=tcblue, fill=tcblue!8, below=1.0cm of topk] (llm) {ReAct Agent\\{\scriptsize (LLM synthesis)}};
    \node[block, draw=tcgreen, fill=tcgreen!8, below=1.0cm of llm] (answer) {Grounded Answer\\+ Sources};
    \node[data, draw=tcgrey!40, fill=tcgrey!4, below=0.4cm of store] (corpus) {8 books, $\sim$3k chunks};

    \draw[arr, color=tcgrey] (query) -- (embed);
    \draw[arr, color=tcgreen] (embed) -- node[above, font=\scriptsize]{384-d vector} (store);
    \draw[arr, color=tccyan] (store) -- node[above, font=\scriptsize]{top 5} (topk);
    \draw[arr, color=tcyellow!80!black] (topk) -- (llm);
    \draw[arr, color=tcgreen] (llm) -- node[right, font=\scriptsize]{+ ensure\_sources()} (answer);
    \draw[arr, dashed, color=tcgrey!40] (corpus) -- (store);

\end{tikzpicture}%
}% end resizebox
\caption[RAG pipeline of the fitness specialist]{RAG pipeline of the fitness specialist. The user query is embedded locally, compared against the normalised chunk vectors via dot product, and the top-$k$ passages are fed to the ReAct agent for synthesis. The \texttt{ensure\_sources()} step guarantees that book citations survive the orchestrator's synthesis.}\label{fig:rag_pipeline}
\end{figure}


\subsection{Telegram: External Server Bridge}\label{sec:data:telegram}

An optional Telegram integration (port 8106) demonstrates the architecture's ability to bridge external stdio-only MCP servers onto the Streamable HTTP bus without forking their code. The upstream \texttt{chigwell/telegram-mcp} project\footnote{\url{https://github.com/chigwell/telegram-mcp}} (116 tools) runs unmodified via \texttt{uv} in an isolated environment; to \texttt{ToolHost} it looks identical to every other server, so a third-party server drops in just like one of our own.

\subsection{Summary}\label{sec:data:summary}

\begin{table}[ht]
\begin{center}
\begin{minipage}{\textwidth}
\caption[Summary of integrated data sources]{Summary of integrated data sources.}\label{tab:data_summary}
\begin{tabularx}{\textwidth}{llcX}
\toprule
\textbf{Source} & \textbf{Server} & \textbf{Tools} & \textbf{Primary Metrics} \\
\midrule
Strava & :8103 & 13 & Activities, load, trends, splits, GPS, records, analysis \\
Garmin & :8104 & 13 & Sleep, HRV, Body Battery, stress, VO\textsubscript{2}max \\
Weather & :8101 & 4 & Forecast, UV, pollen, current conditions \\
Routes & :8102 & 7 & A-B routes, loops, trails, isochrones, elevation \\
Calendar & :8105 & 6 & Events, free slots, event detail, create/update/delete \\
Literature & RAG & 1 & Exercise-science passages (8 books) \\
Telegram & :8106 & 116 & Messages, chats, contacts, media \\
\botrule
\end{tabularx}
\end{minipage}
\end{center}
\end{table}


Tables~\ref{tab:full_tools} and~\ref{tab:full_tools_b} in Appendix~\ref{sec:appendix:tools} expand this summary into the complete tool-by-tool inventory across all seven native MCP servers (the external Telegram bridge's 116 tools are summarised in Table~\ref{tab:data_summary} rather than enumerated individually).
